% Source-aligned English for printed p.220.
\clearpage
\SourcePage{220}

\LayerTarget{tohoku:transition:section-5.7-functorial-formulas}
\LayerTarget{tohoku:equation:5.7.3}
\begin{equation}
\Gamma_\Phi^G=\Gamma_\Psi f_*^G=\Gamma^G\Gamma_\Phi.
\tag{5.7.3}
\end{equation}

\LayerTarget{tohoku:spectral-sequence:5.7-i-supported-equivariant}
Since $f_*^G$ takes injective objects of $C^{X(G)}$ to flabby sheaves, hence
to $\Gamma_\Psi$-acyclic sheaves, the first formula again gives \textit{a
cohomological spectral functor on $C^{X(G)}$, converging to the graded functor
$(H_\Phi^m(X;G,A))$, with initial term}

\LayerTarget{tohoku:equation:5.7.4}
\begin{equation}
\mathrm I_2^{p,q}(A)=H_\Psi^p\bigl(Y,\mathbf H^q(G,A)\bigr).
\tag{5.7.4}
\end{equation}

\LayerTarget{tohoku:obstruction:second-spectral-sequence-with-supports}
The second spectral sequence of Theorem 5.2.1 no longer generalizes as it
stands, since one cannot in general say that, if $B$ is an injective object of
$C^{X(G)}$, then $\Gamma_\Phi(B)$ is a $\Gamma^G$-acyclic $G$-module. But
\textit{suppose that every member of $\Psi$ has a neighborhood belonging to
$\Psi$.}

\LayerTarget{tohoku:construction:extension-by-zero-support-limit}
For every open subset $U\subset Y$, let $A_U=A_{f^{-1}(U)}$ be the abelian
$G$-sheaf on $X$ which is zero on the complement of $f^{-1}(U)$ and agrees
with $A$ on $f^{-1}(U)$. One has
$\Gamma_\Phi(A)=\varinjlim\Gamma_X(A_U)$, where the filtered direct limit is
taken over the ordered filtering set of open subsets $U$ of $Y$ such that
$\overline U\in\Psi$. Since the maps
$\Gamma_X(A_U)\to\Gamma_X(A_V)$ for $U\subset V$ are injective, the functor
$\Gamma^G$ commutes with the filtered direct limit, and one obtains the
functorial identity

\LayerTarget{tohoku:equation:5.7.5}
\begin{equation}
\Gamma_\Phi^G(A)=\varinjlim\Gamma_X^G(A_U),
\tag{5.7.5}
\end{equation}
whence, as in Corollary 1 to Proposition 3.10.1,

\LayerTarget{tohoku:equation:5.7.6}
\begin{equation}
H_\Phi^m(X;G,A)=\varinjlim H^m(X;G,A_U).
\tag{5.7.6}
\end{equation}

\LayerTarget{tohoku:spectral-sequence:ii-a-u}
For every $U$, the group $H^*(X;G,A_U)$ is the abutment of a cohomological
spectral sequence $\mathrm {II}(A,U)$ with initial term

\LayerTarget{tohoku:equation:ii-a-u-initial-term}
\[
\mathrm {II}(A,U)_2^{p,q}=H^p\bigl(G,H^q(X,A_U)\bigr).
\]

\LayerTarget{tohoku:construction:inductive-system-of-spectral-sequences}
For varying $U$, these spectral sequences form an inductive system. It follows
from (5.7.6) that $H_\Phi^*(X;G,A)$ is \textit{the abutment of a
cohomological spectral sequence with initial term}

\LayerTarget{tohoku:equation:5.7.7}
\begin{equation}
\mathrm {II}_2^{p,q}(A)=\varinjlim_U H^p\bigl(G,H^q(X,A_U)\bigr).
\tag{5.7.7}
\end{equation}

\LayerTarget{tohoku:comparison:cohomology-and-filtered-limit}
Of course, this spectral sequence is itself a spectral functor in $A$. In the
most important cases (which we shall make explicit below) one can interchange
$H^p(G,-)$ and $\varinjlim$. Since
$\varinjlim H^q(X,A_U)=H_\Phi^q(X,A)$ (Corollary 1 to Proposition 3.10.1), one
then obtains the usual form of the initial term:

\LayerTarget{tohoku:equation:5.7.7-bis}
\begin{equation}
\mathrm {II}_2^{p,q}=H^p\bigl(G,H_\Phi^q(X,A)\bigr).
\tag{5.7.7 bis}
\end{equation}

\LayerTarget{tohoku:conditions:validity-5.7.7-bis}
\noindent\textit{Conditions for the validity of formula (5.7.7 bis).}

\LayerTarget{tohoku:condition:5.7.7-bis-finite-group}
\noindent a) $G$ is a finite group.

\LayerTarget{tohoku:condition:5.7.7-bis-trivial-action}
\noindent b) $G$ acts trivially on the $H^q(X,A_U)$, and the groups
$H_p(G,\mathbf Z)$ are of finite type.

\LayerTarget{tohoku:condition:5.7.7-bis-essentially-constant}
\noindent c) The inductive system $(H^q(X,A_U))$ is ``essentially constant''
for every $q$; that is, one can find a cofinal system $(U_i)$ such that the
homomorphisms $H^q(X,A_{U_i})\to H^q(X,A_{U_j})$ are isomorphisms.

\noindent\textit{Editorial qualification.} The neighborhood hypothesis above
proves only (5.7.5). The derived filtered-colimit identity (5.7.6) and the
abutment asserted before (5.7.7) additionally require that, for every eligible
$U$, the exact extension-by-zero functor $E_U:A\longmapsto A_U$ take injective
abelian $G$-sheaves to $\Gamma_X^G$-acyclic objects. The comparison with
ordinary cohomology used in (5.7.7 bis) additionally requires $\Phi$ to be
paracompactifying. Without these added hypotheses, no derived filtered-colimit
identity or abutment is asserted here.

\TohokuContinuation{221}
